% !TEX program = lualatex
% ECON 201, Worksheet 4: Specialization and Comparative Advantage. Compile
% with LuaLaTeX so the PDF is tagged (same setup as worksheet03.tex).
% Build from the course root:
%   latexmk -lualatex -cd worksheets/worksheet04.tex && latexmk -c -cd worksheets/worksheet04.tex
\DocumentMetadata{
  pdfstandard=UA-2,
  pdfversion=2.0,
  lang=en-US,
  tagging=on
}
\documentclass[11pt]{article}

\usepackage[letterpaper, margin=0.8in, top=0.5in, bottom=0.45in]{geometry}
\usepackage{fontspec}
\setmainfont{Fira Sans}
\usepackage{amsmath}
\usepackage{amssymb}
\usepackage{graphicx}
\usepackage{booktabs}
\usepackage{array}
% Questions are labelled by hand: under tagging=on, enumitem's enumerate
% did not step its counter, so every item printed the same label.
\newcounter{q}
\newcommand{\Q}[1]{\stepcounter{q}\par\needspace{5\baselineskip}\vspace{3pt}\noindent\hangindent=1.7em\hangafter=1
  \makebox[1.7em][l]{(\alph{q})}#1\par}
\usepackage{xcolor}
\usepackage{needspace}
\definecolor{accent}{HTML}{BF5700}
\usepackage{titlesec}
\titleformat{\section}{\large\bfseries\color{accent}}{}{0pt}{}
\titlespacing*{\section}{0pt}{4pt}{2pt}
\usepackage{hyperref}
\hypersetup{pdftitle={ECON 201 Worksheet: Specialization and Comparative Advantage}, pdfauthor={Div Bhagia}, hidelinks}
\pagestyle{empty}
\setlength{\parindent}{0pt}
\setlength{\parskip}{2pt}

% Ruled answer space: n lines. Plain TeX loop; pgffor's \foreach breaks the
% enumerate counter when used inside an item.
\definecolor{answerrule}{HTML}{B8B8B8}
\newcount\answerlinecount
\newcommand{\answerlines}[1]{%
  \par\answerlinecount=#1
  \loop\ifnum\answerlinecount>0
    \vspace{12pt}\noindent\hspace*{1.7em}%
    {\color{answerrule}\rule{\dimexpr\linewidth-1.7em\relax}{0.4pt}}\par
    \advance\answerlinecount by -1
  \repeat}

% A cell to be filled in. Same width everywhere so the columns line up.
\newcommand{\blank}{\rule{1.15in}{0.3pt}}

\begin{document}

{\LARGE\bfseries\color{accent} Worksheet: Specialization and Comparative Advantage}\hfill{\large ECON 201}

\needspace{7\baselineskip}\section{Definitions from today}

\begingroup\small\setlength{\parskip}{2pt}
\textbf{\color{accent}Opportunity cost.} what you give up by taking one action instead of the next best alternative.\par
\textbf{\color{accent}Absolute advantage.} being able to produce more of a good than someone else, using the same inputs.\par
\textbf{\color{accent}Comparative advantage.} the ability to produce a good at a lower opportunity cost than another producer.\par
\endgroup

\needspace{12\baselineskip}\section{Pablo and Lin's Sunday}\setcounter{q}{0}

Pablo and Lin each spend about 8 hours on Sunday baking bread and making cheese. If they spent the whole day on just one food, they could make:

\begin{center}
\renewcommand{\arraystretch}{1.3}
\begin{tabular}{@{}p{1.2in}>{\centering\arraybackslash}p{1.7in}>{\centering\arraybackslash}p{1.7in}@{}}
\toprule
 & Loaves of bread & Pounds of cheese \\
\midrule
Pablo & 12 & 6 \\
Lin & 2 & 4 \\
\bottomrule
\end{tabular}
\end{center}

They can also split the day between the two foods, and output scales with time. For example, if Pablo spent half the day on bread, he would make 6 loaves, and a quarter of the day on cheese would give him 1.5 pounds.

\Q{Fill in the opportunity cost of each food for each person.}

\begin{center}
\renewcommand{\arraystretch}{1.3}
\begin{tabular}{@{}p{1.2in}p{1.7in}p{1.7in}@{}}
\toprule
 & 1 pound of cheese & 1 loaf of bread \\
\midrule
Pablo & \blank & \blank \\
Lin & \blank & \blank \\
\bottomrule
\end{tabular}
\end{center}

\Q{Who has the absolute advantage in each food?}
\answerlines{1}

\Q{Who has the comparative advantage in each food?}
\answerlines{1}

\Q{Suppose each of them spends 75\% of the day on bread and 25\% on cheese. Fill in what each produces, and the totals.}

\begin{center}
\renewcommand{\arraystretch}{1.3}
\begin{tabular}{@{}p{1.2in}p{1.7in}p{1.7in}@{}}
\toprule
 & Loaves of bread & Pounds of cheese \\
\midrule
Pablo & \blank & \blank \\
Lin & \blank & \blank \\
\midrule
Total & \blank & \blank \\
\bottomrule
\end{tabular}
\end{center}

\Q{Now each makes only the food in which they have the comparative advantage. How much bread and cheese is produced in total?}
\answerlines{1}

\Q{They agree on a deal: Lin sells Pablo 2 pounds of cheese for 2 loaves of bread. What does each of them end up consuming? Is each of them better off than in part (d)?}
\answerlines{2}

\Q{Between which two prices, in loaves of bread per pound of cheese, does trade make both of them better off?}
\answerlines{1}

\needspace{12\baselineskip}\section{One more pair}\setcounter{q}{0}

Ana and Ben work at a campus bakery. In an hour, they can produce:

\begin{center}
\renewcommand{\arraystretch}{1.3}
\begin{tabular}{@{}p{1.2in}>{\centering\arraybackslash}p{1.7in}>{\centering\arraybackslash}p{1.7in}@{}}
\toprule
 & Frosted cupcakes & Pots of coffee \\
\midrule
Ana & 12 & 4 \\
Ben & 4 & 2 \\
\bottomrule
\end{tabular}
\end{center}

\Q{Who has the absolute advantage in each task?}
\answerlines{1}

\Q{Who has the comparative advantage in each task?}
\answerlines{2}

\end{document}
